\section{实证结果分析}

\subsection{福州市老年需求空间分布}

\subsubsection{老年需求强度整体格局}

展示老年需求地图。说明哪些区域老年需求较高。

\subsubsection{高龄照护需求分布}

展示高龄需求或护理需求分布。重点分析是否存在高龄老人集聚片区。

\subsubsection{老年需求热点区识别}

列出需求强度排名前若干网格或街道。

\subsection{社区养老服务供给格局}

\subsubsection{养老服务设施空间分布}

展示养老机构、长者食堂、日间照料中心、养老服务站的整体分布。

\subsubsection{医养协同服务分布}

展示医院、社区卫生服务中心、药店、康复机构等分布。

\subsubsection{代际共享服务分布}

展示公园、社区食堂、公共交通、文化体育空间等分布。

\subsection{养老服务可达性测度结果}

\subsubsection{15 分钟社区养老服务可达性}

展示长者食堂、养老服务站等近距离服务可达性。

\subsubsection{30 分钟照护服务可达性}

展示日间照料中心、社区照护设施可达性。

\subsubsection{代际共享修正前后对比}

展示修正前后的可达性差异。

说明：这一节要突出你们的创新：共享服务确实会改变部分区域的服务评价。

\subsection{供需错配识别结果}

\subsubsection{错配指数空间分布}

展示 $\operatorname{Mismatch}_g$ 地图。说明高错配区集中在哪里。

\subsubsection{四象限分类结果}

展示高需求—低供给、高需求—高供给等四类区域。

\subsubsection{优先补位区识别}

列出高需求—低供给区域，并说明其主要短板。

\subsection{GraphSAGE 高错配风险识别结果}

\subsubsection{模型识别结果}

展示 GraphSAGE 输出的高风险网格图。

\subsubsection{与传统错配指数对比}

比较 GraphSAGE 与基础错配指数结果是否一致。

重点说明：GraphSAGE 能够识别部分受邻近区域影响的潜在风险网格。

\subsubsection{高风险区域特征分析}

总结高风险区域通常具有哪些特征，例如：

\begin{itemize}
\item 老年需求高；
\item 养老服务弱；
\item 医疗协同不足；
\item 公共交通不便；
\item 共享服务支撑不足。
\end{itemize}

\subsection{LightGBM-SHAP 成因解释结果}

\subsubsection{模型拟合效果}

简要说明模型精度。不要占太多篇幅。

\subsubsection{全局特征重要性}

展示 SHAP 特征重要性图。说明哪些变量最影响错配。

\subsubsection{不同短板类型的成因差异}

分析不同区域的主要短板：

\begin{itemize}
\item 助餐不足型；
\item 医养协同不足型；
\item 交通可达性不足型；
\item 共享服务不足型；
\item 综合服务不足型。
\end{itemize}

\subsection{MCLP 服务补位优化结果}

\subsubsection{新增 10 个服务点情景}

展示保守情景下补位点位图。说明覆盖了哪些高需求区域。

\subsubsection{新增 20 个服务点情景}

展示中等情景下补位点位图。说明覆盖效果提升。

\subsubsection{新增 30 个服务点情景}

展示积极情景下补位点位图。说明边际收益是否下降。

\subsubsection{补位前后效果对比}

放一张表：

\begin{table}[H]
\centering
\caption{补位前后效果对比指标}
\begin{tabular}{lllll}
\toprule
\rowcolor{tablehead}
指标 & 补位前 & 新增 10 个 & 新增 20 个 & 新增 30 个 \\
\midrule
覆盖老年人口 &  &  &  &  \\
平均可达时间 &  &  &  &  \\
高错配网格数量 &  &  &  &  \\
服务公平性指数 &  &  &  &  \\
\bottomrule
\end{tabular}
\end{table}
